% === Begin Label Data ===
% ID: 877
% 难度星级: 
% 标签: 
% 备注: 新定义
% 组卷引用次数: 0
% === End  Label Data ===

\begin{problem}{2023}{G}{北京卷}{21}{数列}
数列 $\{a_n\}, \{b_n\}$ 的项数均为 $m(m>2)$, 且 $a_n, b_n \in \{1, 2, \cdots, m\}$, $\{a_n\}, \{b_n\}$ 的前 $n$ 项和分别为 $A_n, B_n$, 并规定 $A_0=B_0=0$. 对于 $k \in \{0, 1, 2, \cdots, m\}$, 定义 $r_k = \max\{i | B_i \leqslant A_k, i \in \{0, 1, 2, \cdots, m\}\}$, 其中, $\max M$ 表示数集 $M$ 中最大的数.

( I ) 若 $a_1=2, a_2=1, a_3=3, b_1=1, b_2=3, b_3=3$, 求 $r_0, r_1, r_2, r_3$ 的值;

( II ) 若 $a_1 \geqslant b_1$, 且 $2r_j \leqslant r_{j+1} + r_{j-1}, j=1, 2, \cdots, m-1$, 求 $r_n$;

( III ) 证明: 存在 $0 \leqslant p < q \leqslant m, 0 \leqslant r < s \leqslant m$, 使得 $A_p + B_s = A_q + B_r$.
\end{problem}

\begin{answer}
( I ) $r_0=0, r_1=1, r_2=1, r_3=2$; ( II ) $r_n=n$.
\end{answer}

\begin{solutions}
解:( I )列表如下, 对比可知 $r_0=0, r_1=1, r_2=1, r_3=2$.

\begin{center}
\begin{tabular}{|c|c|c|c|c|}
\hline
$i$ & $0$ & $1$ & $2$ & $3$ \\
\hline
$a_i$ & & $2$ & $1$ & $3$  \\
\hline
$A_i$ & $0$ & $2$ & $3$ & $6$ \\
\hline
$b_i$ & & $1$ & $3$ & $3$ \\
\hline
$B_i$ & $0$ & $1$ & $4$ & $7$ \\
\hline
$r_k$ & $0$ & $1$ & $1$ & $2$ \\
\hline
\end{tabular}
\end{center}

( II ) 由题可知 $r_0=0$. 因为 $a_1 \geqslant b_1$, 所以 $r_1 \geqslant 1$

因为 $2r_n \leqslant r_{n+1} + r_{n-1}, n=1, 2, \cdots, m-1$,

所以 $r_{n+1} - r_n \geqslant r_n - r_{n-1}$,

所以 $r_n - r_{n-1} \geqslant r_1 - r_0 \geqslant 1, n=1, 2, \cdots, m$,

所以 $r_m = (r_m - r_{m-1}) + (r_{m-1} - r_{m-2}) + \cdots + (r_2 - r_1) + r_1 \geqslant m$ ※

又 $r_m \leqslant m$, (因为题目定义中最大到 $m$)

所以 $r_m=m$.

※中所有的等号均成立, 则 $r_n - r_{n-1} = 1$

综上, $r_n=n, n=0, 1, 2, \cdots, m$.

(3) 对于 $A_k$ 而言, $B_{r_k}$ 是小于等于 $A_k$ 中最大的, 所以 $B_{r_k} \leqslant A_k < B_{r_k+1}$,

即 $B_{r_k} \leqslant A_k < B_{r_k} + b_{r_k+1}$. (注意:为了保证 $b_{r_k+1}$ 存在,需要满足 $r_k < m$)

情况(1)若 $B_m > A_m$, 则 $r_k < m, k=0, 1, \cdots, m$.

于是 $B_{r_k} \leqslant A_k < B_{r_k+1} \Rightarrow B_{r_k} \leqslant A_k < B_{r_k} + b_{r_k+1} \Rightarrow 0 \leqslant A_k - B_{r_k} < b_{r_k+1} \leqslant m \Rightarrow$

$$
\begin{cases}
A_0 - b_0 = 0 \\
0 \leqslant A_1 - B_{r_1} \leqslant m-1 \\
0 \leqslant A_2 - B_{r_2} \leqslant m-1 \\
\cdots \\
0 \leqslant A_m - B_{r_m} \leqslant m-1
\end{cases}
$$

由抽屉原理, $A_0-b_0, A_1-B_{r_1}, \cdots, A_m-B_{r_m}$ 这 $m+1$ 个数必有两个相等,

设为 $A_q - B_{r_q} = A_p - B_{r_p} (p>q)$.

令 $s=r_p, t=r_q$, 则有 $A_p + B_t = A_q + B_s (p>q, s>t)$.

情况(2)若 $B_m = A_m$, 则有 $A_0 + B_m = A_m + B_0$.

情况(3)若 $B_m < A_m$, 则交换集合 A, B 的元素, 按情况(1)处理.
\end{solutions}